\documentclass[11pt]{article}
\usepackage[a4paper,margin=2.2cm]{geometry}
\usepackage{amsmath,amssymb}
\usepackage[T1]{fontenc}
\usepackage{lmodern}
\usepackage[hidelinks]{hyperref}
\setlength{\parindent}{0pt}
\setlength{\parskip}{3pt}
\newcommand{\doi}[1]{\href{https://doi.org/#1}{\texttt{#1}}}
\title{\vspace{-1.2cm}\textbf{Task 0 report --- Paper C (Anu)}\\[2pt]
\large Lamb waves in a magneto-electro-elastic (BaTiO$_3$/CoFe$_2$O$_4$) plate with MGT heat and
NSGT size effect:\\literature sweep, re-anchored problem, and Q1 assessment}
\author{Independent cross-check engineer, for Dr.\ Vipin}
\date{30 October 2026}
\begin{document}
\maketitle
\vspace{-0.5cm}

\section*{1. What this report is}
Condensed from \texttt{paper/literature\_table.md}, \texttt{paper/problem\_v2.md} and
\texttt{paper/q1\_scorecard.md} in \texttt{mee\_anu/}. No new search was run for this report.

\section*{2. What the sweep found}
Fourteen papers were examined. Each ingredient of the plan is individually mature:

\begin{itemize}\small
\item MEE plate waves: Chen, Pan \& Chen, \emph{Int.\ J.\ Solids Struct.} \textbf{44} (2007)
  1073--1085, \doi{10.1016/j.ijsolstr.2006.06.003}; and, with an imperfect interface, Zhang, Shen \&
  Yu, \emph{J.\ Intell.\ Mater.\ Syst.\ Struct.} \textbf{36} (2025) 312--326,
  \doi{10.1177/1045389x241306779}. Nonlocal MEE plates: \emph{J.\ Sound Vib.} \textbf{400} (2017)
  550--563, \doi{10.1016/j.jsv.2017.04.001}.
\item NSGT for MEE plates: \emph{Nanomaterials} \textbf{11} (2021) 87, \doi{10.3390/nano11010087}
  (buckling and vibration, not waves); \emph{Compos.\ Struct.} \textbf{296} (2022) 115878,
  \doi{10.1016/j.compstruct.2022.115878} (thermal, FG, porous).
\item Complex/evanescent Lamb spectra in piezo--piezomagnetic plates: \emph{Materials} \textbf{11}
  (2018) 1186, \doi{10.3390/ma11071186}; FG piezo--piezomagnetic Lamb:
  \emph{Acta Mech.} \textbf{223} (2012) 1081--1091, \doi{10.1007/s00707-012-0612-5}.
\item MGT in plates: \emph{Thin-Walled Struct.} \textbf{180} (2022) 109793,
  \doi{10.1016/j.tws.2022.109793}. ME coefficient vs composition/orientation (the static optimum):
  \emph{AIP Advances} \textbf{7} (2017), \doi{10.1063/1.4994042}.
\end{itemize}

\noindent \textbf{What was not found:} a \emph{thermal} (let alone MGT) MEE \emph{Lamb-wave} plate
paper; and any paper that treats the magnetoelectric optimum as a function of frequency rather than
as a material constant. The first is an open slot but only combination novelty; the second is the
paper.

\begin{table}[h]
\centering\small
\caption{Nearest papers. $\bullet$ treated, $\circ$ partial, --- absent.}
\begin{tabular}{p{5.4cm}p{1.0cm}p{1.0cm}p{1.1cm}p{1.0cm}p{1.0cm}p{1.5cm}}
\hline
Paper & MEE & Plate & $k_i$ & MGT & NSGT & dynamic ME optimum \\
\hline
Chen 2007, \emph{IJSS} 44 & $\bullet$ & $\bullet$ & --- & --- & --- & --- \\
\emph{JIMSS} 36 (2025) & $\bullet$ & $\bullet$ & $\circ$ & --- & --- & --- \\
\emph{JSV} 400 (2017) & $\bullet$ & $\bullet$ & --- & --- & --- & --- \\
\emph{Nanomaterials} 11 (2021) & $\bullet$ & --- & --- & $\circ$ & $\bullet$ & --- \\
\emph{Materials} 11 (2018) & $\circ$ & $\bullet$ & $\bullet$ & --- & --- & --- \\
\emph{AIP Adv.} 7 (2017) & $\bullet$ & --- & --- & --- & --- & --- (static only) \\
\hline
\textbf{this paper} & $\bullet$ & $\bullet$ & $\bullet$ & $\bullet$ & $\bullet$ & $\bullet$ \\
\hline
\end{tabular}
\end{table}

\section*{3. The re-anchored problem}
\textbf{Title.} \emph{The magnetoelectric optimum is not a material constant: a frequency-dependent
design ridge for BaTiO$_3$/CoFe$_2$O$_4$ plates with MGT heat and NSGT.}

\textbf{Finding.} The piezomagnetic volume fraction that maximises the magnetoelectric coupling is
not a material constant: it is set by a mechanical matching condition between the two phases at the
working frequency-thickness, so the optimum traces a ridge in the $(v_f,fd)$ plane whose position
moves with the two NSGT length scales and the MGT relaxation time; the static optimum of the
materials literature is the $fd\to0$ edge of that ridge.

The physics: in a 2-2 composite the ME effect is a \emph{product} property, so strain must be
transferred from the magnetostrictive to the piezoelectric phase; in a dynamic plate that transfer is
efficient only where both phases deform together, which makes the optimum frequency-dependent. The
static limit then provides a rigorous anchor against published values --- a rarity in these three
papers.

\subsection*{Claims}
\begin{enumerate}\small
\item[\textbf{C1}] For each of the four face-condition combinations (electric open/short $\times$
  magnetic open/short) the mode-resolved coefficient $K_{ME}(v_f,fd)$ has an interior maximum, and
  the maximising $v_f^{*}(fd)$ is frequency-dependent. \emph{Fig.\ F4.}
  \emph{Checks:} $v_f=0$ returns the pure piezoelectric plate with $K_{ME}\equiv0$ and $v_f=1$ the
  pure piezomagnetic plate with $K_{ME}\equiv0$, both to $<10^{-6}$.
  \emph{Definitions fixed before Step 2:} $K_E^2$, $K_M^2$ and $K_{ME}$, per mode and per $fd$,
  each with an explicit equation --- the plan's ``K$^2$ from open/short'' was never defined for a
  plate.
\item[\textbf{C2}] The ridge is a mechanical matching condition, anchored by the static limit.
  \emph{Fig.\ F5.} \emph{Checks:} $fd\to0$ reproduces the published static optimum to $<5\%$; two
  phases with identical properties must give $v_f^{*}=0.5$ exactly by symmetry.
\item[\textbf{C3}] NSGT and MGT move the ridge, and the four face conditions give four different
  ridges, hence a design rule. \emph{Fig.\ F7.} \emph{Checks:} $l_{\rm nl}=l_g=0$ returns Step~4;
  $\tau_q\to0$ returns Fourier.
\end{enumerate}

\subsection*{Upgrades}
\textbf{U1 (major):} the dynamic ME ridge and its design rule. \textbf{U2 (minor):} effective-constant
sensitivity across at least two homogenisation schemes, reported as a band --- the direct answer to
the ``constants depend on the source'' objection in \texttt{docs/00\_feasibility\_analysis.md}.
Exceptional points are not used: the MEE plate has no single strong dissipative channel.

\section*{4. Model, boundary conditions, validity}
Fields $u_x,u_z,\varphi,\psi,\theta$ in plane strain, quasi-static electromagnetic approximation
(stated). Constitutive relations
$\sigma=C\varepsilon-eE-qH-\lambda\theta$, $D=e\varepsilon+\varepsilon E+\alpha H+p\theta$,
$B=q\varepsilon+\alpha E+\mu H+m\theta$, with $q$ piezomagnetic, $\alpha$ the magnetoelectric tensor,
$p$ pyroelectric and $m$ pyromagnetic. NSGT is applied to the \emph{mechanical part only} ---
$\alpha$ and the electromagnetic parts are not strain-gradient objects --- and this restriction is
stated, not hidden.

Faces $z=\pm h$: $\sigma_{zz}=\sigma_{xz}=0$, plus the NSGT higher-order traction $P_{izz}=0$
(\texttt{steps/step5\_bcs\_nsgt.md}, approval required before coding); $D_z=0$ or $\varphi=0$;
$B_z=0$ or $\psi=0$; $\partial_z\theta=0$ or $\theta=0$. The 2$\times$2 face-condition table must
appear in every figure caption.

\textbf{Validity window.} $kh\in[0.05,\min(1/\lambda,2h/10\Lambda_{\rm ph})]$ shaded on every figure,
hatched beyond, $h/l_{\max}$ in the caption; plus the MEE-specific bound that the homogenised
constants require the phase size to be much smaller than both the wavelength and the thickness, which
matters across the $v_f$ sweep.

\textbf{Parameters.} Single-phase BaTiO$_3$ and CoFe$_2$O$_4$ sets are to be fixed at Step~1 from
published MEE tables (\texttt{[U]}, not quoted here), with at least two homogenisation schemes;
$v_f=0,0.2,0.4,0.6,0.8,1.0$; $\tau_q\in[10^{-13},10^{-11}]$ s as an assumed range;
$l_{\rm nl},l_g$ as dimensionless $l/h$ sweeps only. The static anchor value must come from a
citable table, not from our own calculation, or C2 cannot be claimed.

\section*{5. Q1 assessment}
\begin{table}[h]
\centering\small
\begin{tabular}{p{5.2cm}p{1.2cm}p{8.2cm}}
\hline
Criterion & Score & Evidence \\
\hline
Novelty vs closest papers & \textbf{2/5} & Every ingredient is individually mature; the dynamic-optimum claim is unpublished but not yet shown to be visible. \\
Physical insight & 3/5 & Strain transfer as a dynamic condition is correct and statically anchored, but it is an observation about known physics. \\
Rigour and verification & 4/5 & Static anchor and the identical-phases degeneracy ($v_f^{*}=0.5$) are machine-checkable; $K_E^2,K_M^2,K_{ME}$ defined before use. \\
Validity of the size-effect model & \textbf{2/5} & Worst host for a gradient theory: NSGT is applied to the homogenised mixture while the real gradients live at phase boundaries the homogenisation has averaged away. \\
Device relevance & 3/5 & A design rule for ME sensors is real, but it is a two-parameter rule, not a device, and design charts already exist in the sensor literature. \\
Reproducibility & 5/5 & Two homogenisation schemes, explicit $K$ definitions, sourced tables, anchor test. \\
\hline
\end{tabular}
\end{table}

\textbf{Total 19/30. Verdict: DOWNGRADE to Q2}, with a Q1 stretch that is not assumed. This matches
the prediction in \texttt{docs/00\_feasibility\_analysis.md} (``lowest Q1 probability as planned'');
the sweep found nothing that lifts it.

\textbf{What would lift it back.} A ridge steep enough to be a design rule: $v_f^{*}(fd)$ must move by
more than $0.15$ in volume fraction across the usable band, and the four face-condition ridges must
differ by more than the ridge width. If those two numbers hold, submit to
\emph{Int.\ J.\ Solids Struct.}; if not, submit to a strong Q2 (\emph{J.\ Intell.\ Mater.\ Syst.\
Struct.} or \emph{Acta Mechanica}) and do not delay.

\textbf{Biggest risk.} \emph{The ME optimum may be flat.} The static materials literature suggests a
broad, shallow $K_{ME}(v_f)$ curve for 2-2 BaTiO$_3$/CoFe$_2$O$_4$; if so, the ridge is noise and the
paper collapses to ``MGT + NSGT for MEE plates''. Test at Step~3 on a reduced (elastic $+$ electric
$+$ magnetic, no heat, no size effect) model.

\textbf{Second risk.} Homogenisation mismatch: two schemes can differ by a factor that moves the ridge
further than the effect being claimed. U2's band is the mitigation, but if the band is wider than the
ridge motion, C1 and C3 both die.

\section*{6. First action requested}
On the reduced model, compute $K_{ME}(v_f)$ at $fd=500$ and $fd=5000$ and report (a) peak height,
(b) peak position, (c) half-width. Three numbers decide Q1 versus Q2.

\section*{7. Corrections to existing documents}
Plan reading item 6 is mis-cited (Sidhardh \& Ray, not Ray). The plan's Step-2 check ``open/short
speed difference gives K$^2$'' must be replaced by three explicit coefficients as in C1. Plan reading
items 3 and 4 (Zakharenko arXiv 2018; Wu et al.\ arXiv 2021) were not resolved to identifiers this
session, so no DOI or URL is claimed for them.

\section*{8. Status}
Task 0 complete; no code written. Journal targets: \emph{Int.\ J.\ Solids Struct.} (Q1, 1.002),
\emph{Eur.\ J.\ Mech.\ A/Solids} (Q1, 1.023); \emph{Wave Motion} is Q2 (0.499, excluded);
\emph{J.\ Intell.\ Mater.\ Syst.\ Struct.} not verified --- do not assume Q1.

\end{document}
